\documentclass[11pt,a4paper]{article}
\usepackage[T1]{fontenc}
\usepackage[utf8]{inputenc}
\usepackage{amsmath,amssymb,amsthm}
\usepackage[margin=1in]{geometry}
\usepackage[colorlinks=true,linkcolor=blue,urlcolor=blue]{hyperref}
\theoremstyle{plain}
\newtheorem{theorem}{Theorem}
\newtheorem{lemma}[theorem]{Lemma}
\newtheorem{corollary}[theorem]{Corollary}
\theoremstyle{definition}
\newtheorem{remark}[theorem]{Remark}
\title{A proof of the conjectured closed form for OEIS A333905}
\author{Adrian Perez Fontelles\\ \small Independent researcher}
\date{25 August 2026}
\begin{document}
\maketitle

\begin{abstract}
OEIS A333905 carries a conjectured closed form for its terms. We prove it. The entry
also posts a generating function $A(x)$; the conjecture is then exactly the statement
that the generating function of the conjectured formula agrees with $A$. The formula is
a $\mathbb{Q}$-linear combination of terms $n^{k}r^{n}$, whose generating function is
computed in closed form by applying $p(\theta)$ to $1/(1-rx)$, with
$\theta=x\,\frac{d}{dx}$. The difference of the two generating functions is then an
explicit rational function; here it collapses to the polynomial $- 2 x^{6} - 2 x^{4} - \frac{x^{3}}{2} - x^{2} + \frac{x}{4}$, of degree
$6$, which proves the closed form for every $n>6$. The computation is exact
rational-function arithmetic, so the conclusion holds for all $n$ at once rather than
for the terms inspected.
\end{abstract}

\noindent\small 2020 Mathematics Subject Classification. 11B37, 05A15, 33F10.\normalsize

\section{The sequence and the conjecture}

OEIS A333905 is ``Lexicographically earliest sequence of distinct positive integers such that a(n) divides the concatenation of a(n+1) to a(n+2)''. It has offset $1$ and begins
\[
a(1),\dots,a(8)\;=\;1, 2, 3, 4, 5, 6, 10, 8,\ \dots
\]
The entry gives the generating function
\begin{equation}\label{eq:gf}
A(x)\;=\;\sum_{n\ge 1} a(n)\,x^{n}\;=\;\frac{x \left(4 x^{7} + 2 x^{5} + x^{4} - x^{2} - 2 x - 1\right)}{2 x^{2} - 1} ,
\end{equation}
and it carries the following comment:
\begin{quote}\small
2\^{}((n-7)/2)*(5 + 2*sqrt(2) + (2*sqrt(2) - 5)*(-1)\^{}n) for n > 6. - \emph{Stefano Spezia}, Oct 23 2021
\end{quote}
As of the ``Last modified'' line on the live entry (2021-10-26, revision 14) the
statement is still recorded as a conjecture, and no proof appears among the entry's links,
formulas or programs.

Write $f(n)$ for the conjectured closed form,
\[
f(n)\;=\;2^{\frac{n}{2} - \frac{7}{2}} \left(\left(-1\right)^{n} \left(-5 + 2 \sqrt{2}\right) + 2 \sqrt{2} + 5\right) .
\]

\section{Closed forms as generating functions}

Let $A(x)=\sum_{n\ge 1}a(n)x^{n}$ and let $\theta=x\,\dfrac{d}{dx}$, so $\theta$
acts on $x^{m}$ as multiplication by $m$.

\begin{lemma}\label{lem:gf}
Let $r\in\mathbb{Q}$ and let $p$ be a polynomial. Then
\[
\sum_{n\ge0}p(n)\,r^{n}x^{n}\;=\;\bigl(p(\theta)\tfrac{1}{1-rx}\bigr)(x),
\]
a rational function of $x$. Consequently the generating function of any
$\mathbb{Q}$-linear combination of terms $n^{k}r^{n}$ is rational and computable in
closed form.
\end{lemma}

\begin{proof}
$\sum_{n\ge0}r^{n}x^{n}=1/(1-rx)$, and applying $\theta$ multiplies the coefficient of
$x^{n}$ by $n$; so $\theta^{k}$ multiplies it by $n^{k}$, and linearity gives the claim.
Starting the sum at an offset $c$ instead of $0$ subtracts the polynomial
$\sum_{n<c}p(n)r^{n}x^{n}$, which changes nothing below.
\end{proof}

\begin{corollary}\label{cor:criterion}
Let $F(x)=\sum_{n\ge 1}f(n)x^{n}$. Then $f(n)=a(n)$ for every $n>d$ if and only
if $F-A$ is a polynomial of degree at most $d$.
\end{corollary}

\begin{proof}
The coefficient of $x^{n}$ in $F-A$ is $f(n)-a(n)$, so that difference vanishes for all
$n>d$ exactly when $F-A$ has no terms above $x^{d}$.
\end{proof}

\section{The computation for A333905}

By Lemma~\ref{lem:gf} the conjectured formula has generating function
\[
F(x)\;=\;\frac{x \left(- 4 x - 5\right)}{4 \left(2 x^{2} - 1\right)} ,
\]
and subtracting the entry's generating function \eqref{eq:gf} gives, after exact
cancellation of rational functions,
\[
F(x)-A(x)\;=\;- 2 x^{6} - 2 x^{4} - \frac{x^{3}}{2} - x^{2} + \frac{x}{4} ,
\]
a polynomial of degree $6$.

\begin{theorem}
For every $n>6$,
\[
a(n)\;=\;2^{\frac{n}{2} - \frac{7}{2}} \left(\left(-1\right)^{n} \left(-5 + 2 \sqrt{2}\right) + 2 \sqrt{2} + 5\right) .
\]
\end{theorem}

\begin{proof}
Immediate from Corollary~\ref{cor:criterion} and the displayed value of $B$.
\end{proof}

\begin{remark}
The polynomial $F-A$ is not an error term: it records the boundary of the claim. A closed
form of this kind typically fails at a few initial indices, where the sequence's own
definition has not yet settled into its eventual pattern, and $F-A$ collects exactly
those discrepancies. Above that range the identity is exact.
\end{remark}

\section{Verification}

Every number above is an output of a script written separately from this note, and two
independent checks were run.

First, the generating function was checked against the entry rather than assumed: the
Taylor coefficients of \eqref{eq:gf} were expanded and compared term by term with the DATA
section of A333905, agreeing on all 13 terms compared.

Second, the closed form was evaluated directly on the published terms, without reference
to any generating function: $f(n)$ was computed in exact rational arithmetic and compared
with the entry's own value for every $n>6$ in range, agreeing in all 45 cases
from $n=7$ upward. This is what Corollary~\ref{cor:criterion} predicts from the
degree of $F-A$, and it is a genuinely different computation from the symbolic one.

The symbolic step is exact throughout. The difference is obtained by rational-function
cancellation in $K$, not by series truncation, so the identity $B=- 2 x^{6} - 2 x^{4} - \frac{x^{3}}{2} - x^{2} + \frac{x}{4}$ is an
identity of functions and the theorem holds for all $n>6$ simultaneously.

\begin{thebibliography}{9}
\bibitem{oeis} The OEIS Foundation, \emph{The On-Line Encyclopedia of Integer Sequences},
\url{https://oeis.org}, 2026. Sequence A333905.
\bibitem{kp} M.~Kauers and P.~Paule, \emph{The Concrete Tetrahedron}, Springer, 2011,
Chapter 7 (holonomic closure properties and the translation between recurrences and
differential equations).
\bibitem{stanley} R.~P.~Stanley, \emph{Enumerative Combinatorics}, Volume 2, Cambridge
University Press, 1999, Chapter 6 (algebraic and D-finite generating functions).
\end{thebibliography}
\end{document}
